\section*{الفصل \ref{ch:b1:vspaces} --- الفضاءات المتجهية}

\begin{solution}{exo:b1:vspaces:1}
\begin{enumerate}
  \item نعم: فهي تحتوي $0$، والمعادلة المعرِّفة خطية
        (ومستقرة بالمقدار $x + \lambda y$).
  \item لا: فهي لا تحتوي $(0,0,0)$.
  \item لا: فالمتجهتان $(1, 0)$ و $(0, -1)$ تنتميان (لأن $xy = 0$)، ومجموعهما $(1,
        -1)$ لا ينتمي (لأن $xy = -1 < 0$).
  \item نعم: فالمقدار $0$ ينعدم عند $1$؛ و $(P + \lambda Q)(1) = P(1) +
        \lambda Q(1) = 0$.
  \item نعم: فالدالة المعدومة محدودة؛ وإذا كان $\abs f \leq M$ و
        $\abs g \leq M'$، فإن $\abs{f + \lambda g} \leq M +
        \abs\lambda M'$.
\end{enumerate}
\end{solution}

\begin{solution}{exo:b1:vspaces:2}
$(1,2,1) = a(1,0,1) + b(1,1,0)$ يقتضي $a + b = 1$ و $b = 2$ و $a =
1$: وهذا غير متوافق (لأن $a + b = 3 \neq 1$): فليست في \omterm{def:b1:vspaces:span}{الفضاء المولَّد}. و
$(2,1,1) =
a(1,0,1) + b(1,1,0)$: $b = 1$ و $a = 1$ و $a + b = 2$: وهذا متسق، ومنه
$(2,1,1) = (1,0,1) + (1,1,0)$، فهي في \omterm{def:b1:vspaces:span}{الفضاء المولَّد}.

والمعادلة: $(x, y, z) = (a + b, b, a)$ يعني $x = y + z$: ومنه \omterm{def:b1:vspaces:span}{فالفضاء المولَّد} هو
المستوي $\{x - y - z = 0\}$.
\end{solution}

\begin{solution}{exo:b1:vspaces:3}
العائلة الأولى: يعطي $\lambda(1,1,0) + \mu(1,0,1) + \nu(0,1,1) = 0$
أن $\lambda + \mu = 0$ و $\lambda + \nu = 0$ و $\mu + \nu = 0$: وبالجمع،
$2(\lambda + \mu + \nu) = 0$، وبطرح كل معادلة،
$\lambda = \mu = \nu = 0$: فهي \omterm{def:b1:vspaces:free}{حرة}.

والثانية: $(2,4,6) = 2(1,2,3)$: فهي مرتبطة.

والثالثة: أربع متجهات في $\R^3$ --- وهي مرتبطة بالضرورة ما إن يتوفر
البُعد (\cref{ch:b1:findim})؛ ومباشرةً: $(0,1,1) = -(1,0,0) +
0\cdot(1,1,0) + (1,1,1)$، وفعلًا $(-1,0,0) + (1,1,1) = (0,1,1)$: أي
علاقة غير بديهية.
\end{solution}

\begin{solution}{exo:b1:vspaces:4}
لكثيرات الحدود $1, (X-1), (X-1)^2, (X-1)^3$ درجات متمايزة $0,
1, 2, 3$: فهي \omterm{def:b1:vspaces:free}{حرة} (\cref{prop:b1:vspaces:freecriteria}~(1))؛ وأربع
متجهات \omterm{def:b1:vspaces:free}{حرة} مولِّدة (لأن كل $P \in \R_3[X]$ ينتشر بقوى $X -
1$، مثلًا بتايلور من أجل كثيرات الحدود، قارن برهان
\cref{prop:b1:poly:multiplicity}): أي \omterm{def:b1:vspaces:free}{أساس}. ومن أجل $X^3$، تايلور عند
$1$: $P = X^3$ و $P(1) = 1$ و $P'(1) = 3$ و $P''(1) = 6$ و $P'''(1) =
6$:
\[
X^3 = 1 + 3(X - 1) + 3(X-1)^2 + (X-1)^3 ,
\]
\omterm{prop:b1:vspaces:coordinates}{فالإحداثيات} $(1, 3, 3, 1)$ (وهو سطر باسكال، كما هو متوقع من $X^3 =
((X-1)+1)^3$).
\end{solution}

\begin{solution}{exo:b1:vspaces:5}
$F \cap G$: يعطي الشرطان $x = y = z$ و $x = t = 0$ معًا
$x = 0$، ومنه $y = z = 0$ و $t = 0$: ومنه فالتقاطع هو
$\{0\}$. والمجموع: من أجل $(x,y,z,t)$ معطاة، ابحث عن $(a,a,a,b) \in F$ و
$(0,c,d,0) \in G$ مجموعهما هي: $a = x$ و $b = t$ و $c = y - x$ و $d =
z - x$: وهذا ممكن دائمًا. ومنه $\R^4 = F \oplus G$، و
\[
(1,2,3,4) = (1,1,1,4) + (0,1,2,0) .
\]
\end{solution}

\begin{solution}{exo:b1:vspaces:6}
كل عنصر من $G$ هو $\lambda u$؛ فإذا تقارب، فإن (لأن
$\lambda u_n = \lambda(-1)^n$ له نهايتا متتاليتين جزئيتين $\pm
\lambda$) يكون بالضرورة $\lambda = 0$: أي $F \cap G = \{0\}$.

و $F + G$ ليس كل شيء: فهو يتألف من متتاليات من الشكل
$c_n + \lambda(-1)^n$ مع $(c_n)$ متقاربة. والمتتالية $v_n =
n$ ليست من هذا الشكل (لأن $v_n - \lambda(-1)^n$ غير محدودة، فهي غير
متقاربة أبدًا). ومنه $F \oplus G \subsetneq$ (فضاء كل المتتاليات).
\end{solution}

\begin{solution}{exo:b1:vspaces:7}
($\supseteq$) يقع كلٌّ من $F \cap G$ و $F \cap H$ في $F$، ويقع
مجموعهما في $G + (F \cap H)$: ويتبع الاحتواء لأن الطرف الأيسر
\omterm{def:b1:vspaces:subspace}{فضاء جزئي} يحتوي القطعتين --- وملموسًا، فالعنصر $g +
h$ مع $g \in F\cap G$ و $h \in F \cap H$ يقع في $F$ (مجموع عنصرين
من $F$) وفي $G + (F \cap H)$.

($\subseteq$) ليكن $x \in F$ مع $x = g + h$ و $g \in G$ و $h \in F
\cap H$. عندئذ $g = x - h \in F$ (فرق عنصرين من $F$)، ومنه
$g \in F \cap G$، و $x = g + h \in (F \cap G) + (F \cap H)$.

ومثال مضاد للتوزيعية الكاملة في $\R^2$: $F =
\operatorname{Vect}(1,1)$ و $G = \operatorname{Vect}(1,0)$ و $H =
\operatorname{Vect}(0,1)$. عندئذ $G + H = \R^2$، ومنه $F \cap (G+H) =
F$، بينما $F \cap G = F \cap H = \{0\}$: فالطرف الأيمن هو $\{0\}
\neq F$.
\end{solution}

\begin{solution}{exo:b1:vspaces:8}
\begin{enumerate}
  \item افترض $\sum_{i} \lambda_i \eu^{a_i x} = 0$ من أجل كل $x$.
        اضرب في $\eu^{-a_p x}$: $\lambda_p + \sum_{i < p}
        \lambda_i \eu^{(a_i - a_p)x} \to \lambda_p$ عندما $x \to
        +\infty$ (لأن كل أُسّ $a_i - a_p < 0$). والطرف الأيسر يساوي
        $0$ تمامًا، ومنه $\lambda_p = 0$؛ وكرّر نزولًا.
  \item ليكن $a\cos x + b \sin x + c \cos 2x + d\sin 2x = 0$ من أجل كل
        $x$. قوّم عند $x = 0$: $a + c = 0$؛ وعند $x = \pi$: $-a + c
        = 0$؛ ومنه $a = c = 0$، وتُردّ العلاقة إلى $b\sin
        x + d \sin 2x = 0$. وقوّم عند $x = \frac\pi2$: $b = 0$؛
        ثم عند $x = \frac\pi4$: $d = 0$.
  \item افترض $\sum \lambda_i \abs{x - a_i} = 0$ من أجل كل $x$. والدالة
        $\sum_{i \neq j} \lambda_i\abs{x - a_i}$ \omterm{def:b1:derivative:def}{قابلة للاشتقاق}
        عند $a_j$ (لأن كل حدّ كذلك، بعيدًا عن ركنه
        الخاص)، ومنه فعلى $-\lambda_j \abs{x - a_j}$، وهي فرقهما،
        أن تكون \omterm{def:b1:derivative:def}{قابلة للاشتقاق} عند $a_j$ كذلك --- وهذا يفرض
        $\lambda_j = 0$ (لأن للمقدار $\abs{\,\cdot\,}$ ركنًا). ويصحّ هذا
        من أجل كل $j$.
\end{enumerate}
\end{solution}

\begin{solution}{exo:b1:vspaces:9}
ليكن $h \in H$. ولأن $h \in H \subseteq F + H = F + G$، اكتب $h = f
+ g$ مع $f \in F$ و $g \in G$. عندئذ $f = h - g \in H$ (لأن الحدّين
في $H$، باستعمال $G \subseteq H$)، ومنه $f \in F \cap H = F \cap G
\subseteq G$، و $h = f + g \in G$. ومنه $H \subseteq G$، ومع
الفرض $G \subseteq H$: تقع المساواة.

ومثال مضاد دون $G \subseteq H$: في $\R^2$، خذ $F =
\operatorname{Vect}(1,0)$ و $G = \operatorname{Vect}(0,1)$ و $H =
\operatorname{Vect}(1,1)$: عندئذ $F + G = F + H = \R^2$ و $F \cap G
= F \cap H = \{0\}$، ومع ذلك $G \neq H$.
\end{solution}

\begin{solution}{exo:b1:vspaces:10}
المجموعتان فضاءان جزئيان (لأن الشروط المعرِّفة خطية و
تتحقق من أجل $0$). والتفكيك: من أجل $P \in \R[X]$،
\[
P(X) = \underbrace{\frac{P(X) + P(-X)}{2}}_{\in\,\mathcal P}
+ \underbrace{\frac{P(X) - P(-X)}{2}}_{\in\,\mathcal I},
\]
\omterm{def:b1:poly:def}{وكثير الحدود} الزوجي والفردي معًا يحقق $P = -P$، ومنه $P = 0$:
\omterm{def:b1:vspaces:sum}{فالمجموع مباشر} ويساوي $\R[X]$.

وليكن الآن $P = \sum_k a_k X^k$ زوجيًا. عندئذ يكون $P(X) - P(-X) = 2
\sum_{k \text{ فردي}} a_k X^k$ هو \omterm{def:b1:poly:def}{كثير الحدود} المعدوم، ومنه ينعدم كل
معامل ذي درجة فردية (\cref{def:b1:poly:def}): $P \in
\operatorname{Vect}(1, X^2, X^4, \dots)$، أي $P = Q(X^2)$ من أجل
\omterm{def:b1:poly:def}{كثير حدود} $Q$. وبالعكس كل \omterm{def:b1:poly:def}{كثير حدود} في $X^2$ زوجي.
\end{solution}

\begin{solution}{exo:b1:vspaces:11}
افترض $a(x_1 + x_2) + b(x_2 + x_3) + c(x_3 + x_1) = 0$.
وبإعادة التجميع على \omterm{def:b1:vspaces:free}{العائلة الحرة} $(x_1, x_2, x_3)$:
\[
(a + c)\,x_1 + (a + b)\,x_2 + (b + c)\,x_3 = 0
\implies a + c = a + b = b + c = 0 .
\]
وبطرح المعادلتين الأوليين نجد $c = b$؛ ثم تعطي الثالثة
$2b = 0$، ومنه $b = c = 0$، ثم $a = 0$: \omterm{def:b1:vspaces:free}{فالعائلة حرة}.

وأمّا من أجل أربع متجهات فالعائلة النظيرة مرتبطة \emph{دائمًا}:
\[
(x_1 + x_2) - (x_2 + x_3) + (x_3 + x_4) - (x_4 + x_1) = 0
\]
تركيبةٌ معدومة غير بديهية (بالمعاملات $1, -1, 1, -1$)،
مهما كانت $(x_1, x_2, x_3, x_4)$. فزوجية طول الدورة
هي التي تفصل.
\end{solution}

\begin{solution}{exo:b1:vspaces:12}
\begin{enumerate}
  \item إذا كان $F_1 \subseteq F_2$ أو $F_2 \subseteq F_1$، فالاتحاد
        هو أحدهما، ومنه فهو قطعي. وإلا فاختر $x \in F_1
        \setminus F_2$ و $y \in F_2 \setminus F_1$، وانظر في
        $x + y$. فإذا كان $x + y \in F_1$، فإن $y = (x + y) - x \in
        F_1$: وهذا تناقض. وإذا كان $x + y \in F_2$، فإن $x \in F_2$:
        وهذا تناقض. ومنه $x + y \notin F_1 \cup F_2$، و $E \neq
        F_1 \cup F_2$.
  \item افترض بالخلف أن $E = F_1 \cup \dots \cup
        F_k$، مع اختيار $k$ \emph{أصغريًا} بين كل هذه
        التغطيات. وتمنع الأصغرية $F_1 \subseteq F_2 \cup
        \dots \cup F_k$ (وإلا فأسقط $F_1$)، ومنه يوجد $x \in F_1$
        مع $x \notin F_i$ من أجل كل $i \geq 2$. ولأن $F_1$
        قطعي، اختر $y \notin F_1$. ومن أجل كل سلّم $t$، تقع
        المتجهة $y + t x$ في $F_i$ ما. وهي لا تقع أبدًا في
        $F_1$: وإلا لكان $y = (y + tx) - tx \in F_1$ (لأن $x \in
        F_1$). \omterm{def:b1:structures:field}{والحقل} غير منته، فاختر $k$ سلّمًا
        متمايزًا $t_1, \dots, t_k$: فتقع المتجهات $k$ الممثَّلة بالمقدار $y + t_j x$
        في الفضاءات الجزئية $k - 1$ وهي $F_2, \dots, F_k$، وتقع اثنتان
        منها، ولنقل $y + t x$ و $y + t' x$ مع $t \neq t'$،
        في $F_i$ نفسه ($i \geq 2$). عندئذ فرقهما
        $(t - t')x \in F_i$، ومنه $x \in F_i$: وهذا تناقض. ومنه
        فلا وجود لتغطية منتهية بفضاءات جزئية قطعية.
\end{enumerate}
\end{solution}

\begin{solution}{pb:b1:vspaces:1}
\textbf{1.} المقدار $L_i$ جداءُ $n$ عاملًا خطيًا مقسومًا على
ثابت غير معدوم (لأن المقادير $x_i$ متمايزة)، ومنه $\deg L_i = n$.
وبالتقويم عند $x_j$ مع $j \neq i$: ينعدم العامل $X - x_j$ في
البسط، ومنه $L_i(x_j) = 0$. وعند $x_i$، يتطابق البسط و
المقام: $L_i(x_i) = 1$.

\textbf{2.} افترض $\sum_i \lambda_i L_i = 0$. وقوّم عند $x_j$:
فتموت كل الحدود إلا $\lambda_j L_j(x_j) = \lambda_j$، ومنه
$\lambda_j = 0$ من أجل كل $j$: \omterm{def:b1:vspaces:free}{فالعائلة حرة}.

\textbf{3.} لتكن $D = P - \sum_i P(x_i) L_i$. عندئذ $\deg D \leq n$
و، حسب السؤال 1، $D(x_j) = P(x_j) - P(x_j) = 0$ من أجل النقاط المتمايزة $n + 1$
$x_0, \dots, x_n$. ولكثير حدود غير معدوم من درجة
$\leq n$ عددُ $n$ جذرًا على الأكثر (\cref{cor:b1:poly:nroots})، ومنه $D
= 0$. ومنه فكل $P \in \R_n[X]$ تركيبةٌ من المقادير $L_i$:
فالعائلة مولِّدة، ومع السؤال 2، \omterm{def:b1:vspaces:free}{أساس}.

\textbf{4.} من أجل $y_0, \dots, y_n$ معطاة، يكون \omterm{def:b1:poly:def}{لكثير الحدود} $P = \sum_i
y_i L_i$ درجةٌ $\leq n$ وهو يقايس. والوحدانية: \omterm{def:b1:poly:def}{لكثير حدود}
مقايِس $P$، حسب السؤال 3، \omterm{prop:b1:vspaces:coordinates}{إحداثيات} $(P(x_0), \dots,
P(x_n)) = (y_0, \dots, y_n)$ في \omterm{def:b1:vspaces:free}{الأساس} $(L_i)$، \omterm{prop:b1:vspaces:coordinates}{والإحداثيات}
في \omterm{def:b1:vspaces:free}{أساس} وحيدة (\cref{prop:b1:vspaces:coordinates}). و
\omterm{prop:b1:vspaces:coordinates}{إحداثيات} $P$ في \omterm{def:b1:vspaces:free}{أساس} لاغرانج هي \emph{قيمه عند
العقد} --- وهذا كل مقصود هذا \omterm{def:b1:vspaces:free}{الأساس}.

\textbf{5.} طبّق السؤال 3 على $P = X^k$ ($0 \leq k \leq n$):
\[
X^k = \sum_{i=0}^{n} x_i^{k} L_i ,
\]
ويعطي $k = 0$ أن $\sum_i L_i = 1$.

\textbf{6.} $\deg N_k = k$ بالضبط: فالعائلة $(N_0, \dots, N_n)$
سلّمُ درجات في $\R_n[X]$، ومنه فهي \omterm{def:b1:vspaces:free}{أساس} حسب
\cref{ex:b1:vspaces:staircase} (والحرية من
\cref{prop:b1:vspaces:freecriteria}~(1)، والتوليد بنزول منته
على الدرجة).

\textbf{7.} من العلاقة التراجعية،
\[
f[x_0, x_1] = \frac{f(x_1) - f(x_0)}{x_1 - x_0},
\qquad
f[x_0, x_1, x_2] = \frac{f[x_1, x_2] - f[x_0, x_1]}{x_2 - x_0}.
\]
ومن أجل $f(x) = x^2$:
\[
f[x_0, x_1] = \frac{x_1^2 - x_0^2}{x_1 - x_0} = x_0 + x_1,
\]
ثم
\[
f[x_0, x_1, x_2]
= \frac{(x_1 + x_2) - (x_0 + x_1)}{x_2 - x_0}
= \frac{x_2 - x_0}{x_2 - x_0} = 1 .
\]

\textbf{8.} $\deg S \leq n$ لأن درجة $Q, R$ هي $\leq n - 1$.
وعند $x_0$: $S(x_0) = \frac{-(x_0 - x_n) R(x_0)}{x_n - x_0} =
R(x_0) = f(x_0)$. وعند $x_n$: $S(x_n) = \frac{(x_n - x_0)
Q(x_n)}{x_n - x_0} = Q(x_n) = f(x_n)$. وعند عقدة داخلية $x_i$
($1 \leq i \leq n-1$)، تأخذ $Q$ و $R$ معًا القيمة $f(x_i)$،
ومنه
\[
S(x_i) = \frac{(x_i - x_0) - (x_i - x_n)}{x_n - x_0}\, f(x_i)
= f(x_i) .
\]

\textbf{9.} بالاستقراء على عدد العقد. عقدة واحدة: فالمقايِس
هو الثابت $f(x_0) = f[x_0]$. افترض الادعاء
من أجل $k$ عقدة ولتقايس $S$ عند $x_0, \dots, x_k$؛ فبالوحدانية
(السؤال 4)، تُعطى $S$ بمبرهنة أيتكن المساعدة من $R$
(بالعقد $x_0, \dots, x_{k-1}$) و $Q$ (بالعقد $x_1, \dots, x_k$).
ومعامل $X^{k}$ في $S$ هو
\[
\frac{[X^{k-1}]\,Q - [X^{k-1}]\,R}{x_k - x_0}
= \frac{f[x_1, \dots, x_k] - f[x_0, \dots, x_{k-1}]}{x_k - x_0}
= f[x_0, \dots, x_k]
\]
حسب فرض الاستقراء والعلاقة التراجعية المعرِّفة.

\textbf{10.} ليقايس $P_k$ الدالةَ $f$ عند $x_0, \dots, x_k$. فللفرق
$P_k - P_{k-1}$ درجةٌ $\leq k$ وهو ينعدم عند
$x_0, \dots, x_{k-1}$، ومنه، بمبرهنة العامل مطبَّقةً $k$ مرة
(\cref{thm:b1:poly:factor})، يساوي $c\,N_k$ من أجل ثابت $c$؛
وبمقارنة معاملات $X^{k}$ وباستعمال السؤال 9، $c =
f[x_0, \dots, x_k]$. والتلسكوب انطلاقًا من $P_0 = f(x_0) N_0$ يعطي
صيغة نيوتن. وأمّا الصورة \omterm{def:b1:topology:closed}{المغلقة}، فاكتب $P_k = \sum_{i \leq
k} f(x_i) L_i$ (بلاغرانج، على العقد $x_0, \dots, x_k$) واقرأ
معامل $X^{k}$: فيسهم كل $L_i$ بالمقدار
$\frac{1}{\prod_{j \neq i}(x_i - x_j)}$، ومنه
\[
f[x_0, \dots, x_k] = \sum_{i=0}^{k}
\frac{f(x_i)}{\prod_{j \neq i,\, j \leq k}(x_i - x_j)} .
\]
والطرف الأيمن ثابت تحت أيّ \omterm{def:b1:counting:objects}{تبديلة} للعقد، ومنه
فالفرق المقسوم لا يتعلق بترتيبها.

\textbf{11.} إذا كان $P = a X^m + (\text{درجات أدنى})$، فإن مبرهنة
ثنائي الحدّ تعطي
\[
\Delta P = a\bigl((X+1)^m - X^m\bigr) + \dots
= a\,m\,X^{m-1} + (\text{درجات أدنى}),
\]
لأن $(X+1)^m - X^m = m X^{m-1} + \dots$ ولأن الجزء ذا الدرجات الأدنى
من $P$ يسهم بدرجة $\leq m - 2$ بعد $\Delta$ (أو
بحدود من درجة $\leq m-2$). ومنه $\deg \Delta P = m - 1$ بمعامل
مهيمن $m a$. وأمّا الثابت $c$ فيعطي $\Delta c = c - c = 0$.

\textbf{12.} $\deg B_k = k$: أي سلّم، ومنه فهي \omterm{def:b1:vspaces:free}{أساس} للمقدار
$\R_n[X]$ (\cref{ex:b1:vspaces:staircase}). ومن أجل $\Delta B_k$
($k \geq 1$)، عمّل الجداء المشترك:
\begin{align*}
k!\,\Delta B_k
&= (X+1)X\cdots(X-k+2) - X(X-1)\cdots(X-k+1) \\
&= X(X-1)\cdots(X-k+2)\,\bigl[(X+1) - (X-k+1)\bigr] \\
&= k\,X(X-1)\cdots(X-k+2),
\end{align*}
ومنه $\Delta B_k = \frac{X(X-1)\cdots(X-k+2)}{(k-1)!} = B_{k-1}$.

\textbf{13.} اكتب $P = \sum_{k=0}^{n} c_k B_k$ (\omterm{def:b1:vspaces:free}{أساس}، السؤال
12). وطبّق $\Delta^{j}$: فحسب السؤال 12، $\Delta^{j} P = \sum_{k
\geq j} c_k B_{k-j}$. وقوّم عند $0$: $B_0(0) = 1$ و $B_m(0) =
0$ من أجل $m \geq 1$ (لأن العامل $X$ ينعدم)، ومنه
$\bigl(\Delta^{j}P\bigr)(0) = c_j$. وهذه هي صيغة الفروق
الأمامية.

\textbf{14.} بالاستقراء على $k$. من أجل $k = 0$ تُقرأ المتطابقة
$P(0) = P(0)$. افترضها من أجل $k$ وطبّقها على $\Delta P$:
\[
\bigl(\Delta^{k+1} P\bigr)(0)
= \sum_{j=0}^{k} (-1)^{k-j}\binom kj \bigl(P(j+1) - P(j)\bigr).
\]
واجمع معامل $P(i)$: فهو $(-1)^{k-i+1}\binom
k{i-1}\cdot(-1)^{0}$ من المجموع الأول (مزاحًا) و
$-(-1)^{k-i}\binom ki$ من الثاني --- ومعًا
\[
(-1)^{k+1-i}\Bigl(\binom k{i-1} + \binom ki\Bigr)
= (-1)^{k+1-i}\binom{k+1}i
\]
بقاعدة باسكال، وهذه هي المتطابقة عند الرتبة $k + 1$.

\textbf{15.} بتكرار السؤال 11 انطلاقًا من الدرجة $n$ بمعامل
مهيمن $a_n$: بعد $\Delta$ واحد، تكون الدرجة $n-1$ والمعامل
المهيمن $n a_n$؛ وبعد اثنين، $n(n-1)a_n$؛ وبعد $n$ خطوة،
تكون الدرجة $0$ والقيمة $n(n-1)\cdots 1\, a_n = n!\,a_n$، وهي ثابت.
ويقتله $\Delta$ آخر: $\Delta^{n+1}P = 0$.

\textbf{16.} إذا كان $m \geq k$: $B_k(m) = \binom mk \in \N$. وإذا كان $0
\leq m < k$: فأحد عوامل $m(m-1)\cdots(m-k+1)$ معدوم، ومنه
$B_k(m) = 0$. وإذا كان $m = -q$ مع $q \geq 1$:
\[
B_k(-q) = \frac{(-q)(-q-1)\cdots(-q-k+1)}{k!}
= (-1)^k\,\frac{q(q+1)\cdots(q+k-1)}{k!}
= (-1)^k \binom{q+k-1}{k},
\]
وهو عدد صحيح. ومنه فكل $B_k$ يرسل $\Z$ داخل $\Z$.

\textbf{17.} ($\Leftarrow$) إذا كان $P = \sum_k c_k B_k$ مع $c_k \in
\Z$، فإنه من أجل $m \in \Z$، $P(m) = \sum_k c_k B_k(m) \in \Z$ حسب
السؤال 16. ($\Rightarrow$) إذا كان $P$ صحيح القيم، فإن
إحداثياته هي $c_k = \bigl(\Delta^k P\bigr)(0) = \sum_{j=0}^k
(-1)^{k-j}\binom kj P(j)$ (السؤالان 13 و 14)، وهي تركيبة
صحيحة من الأعداد الصحيحة $P(0), \dots, P(k)$. وهذا هو
تمييز بوليا لكثيرات الحدود صحيحة القيم.

\textbf{18.} ضع $Q(X) = P(X + a)$، وهو \omterm{def:b1:poly:def}{كثير حدود} من الدرجة $\leq
n$ مع $Q(0), Q(1), \dots, Q(n) \in \Z$. وإحداثياته في
$(B_k)_{k \leq n}$ هي $c_k = \sum_{j \leq k}(-1)^{k-j}\binom kj
Q(j) \in \Z$ (لأن السؤال 14 لا يستعمل إلا القيم عند $0, \dots, k
\leq n$). وحسب السؤال 17 ($\Leftarrow$)، يكون $Q$ صحيح القيم على
$\Z$ كله، ومنه فكذلك $P(X) = Q(X - a)$.

\textbf{19.} جداء $k$ عددًا صحيحًا متتاليًا هو $m(m-1)
\cdots(m-k+1) = k!\,B_k(m)$ من أجل $m \in \Z$ ما، و $B_k(m) \in
\Z$ حسب السؤال 16: ومنه فالجداء يقبل القسمة على $k!$.

\textbf{20.} قيم $P = \frac{X(X+1)(2X+1)}{6}$ عند $0,1,2,3$:
$0, 1, 5, 14$. وجدول الفروق: سطر $\Delta$ هو $1, 4, 9$؛
وسطر $\Delta^2$ هو $3, 5$؛ وسطر $\Delta^3$ هو $2$. ومنه، حسب السؤال 13،
\[
P = 0\cdot B_0 + 1\cdot B_1 + 3\,B_2 + 2\,B_3 ,
\]
\omterm{prop:b1:vspaces:coordinates}{بإحداثيات} صحيحة: فالمقدار $P$ صحيح القيم (السؤال 17)،
بينما ليست معاملاته الوحيدة الحدّ $\frac13, \frac12, \frac16$
أعدادًا صحيحة. وبالحساب المباشر:
\begin{align*}
\Delta P &= \frac{(X+1)(X+2)(2X+3) - X(X+1)(2X+1)}{6} \\
&= \frac{(X+1)\bigl[(X+2)(2X+3) - X(2X+1)\bigr]}{6}
= \frac{(X+1)(6X+6)}{6} = (X+1)^2 .
\end{align*}
وبتلسكوب $P(m) = \sum_{j=0}^{m-1}\Delta P(j) = \sum_{j=1}^{m}
j^2$ (مع $P(0) = 0$): تنتج صيغة مجموع المربعات.

\textbf{21.} للقيم $2^i$ عند $i = 0, \dots, n$ جدولُ فروق
ثابت يساوي $1$ على \omterm{def:b1:topology:closure}{الحافة} اليسرى: فالمقدار $\Delta^k$ للمتتالية
$(2^i)$ هو $(2^i)$ مرة أخرى (لأن $2^{i+1} - 2^i = 2^i$)، ومنه
$\bigl(\Delta^k P\bigr)(0) = 2^0 = 1$ من أجل كل $k \leq n$، و $P =
B_0 + B_1 + \dots + B_n$ حسب السؤال 13. عندئذ
\[
P(n+1) = \sum_{k=0}^{n}\binom{n+1}{k}
= 2^{n+1} - \binom{n+1}{n+1} = 2^{n+1} - 1 \neq 2^{n+1}:
\]
فينكسر النمط عند أول نقطة غير مضبوطة.

\textbf{22.} حسب السؤال 12، $B_k = \Delta B_{k+1}$، ومنه
\[
\sum_{j=0}^{m-1} B_k(j)
= \sum_{j=0}^{m-1}\bigl(B_{k+1}(j+1) - B_{k+1}(j)\bigr)
= B_{k+1}(m) - B_{k+1}(0) = B_{k+1}(m).
\]
ومن أجل $j < k$ تنعدم الحدود $B_k(j)$، ومنه يبدأ المجموع فعلًا عند
$j = k$: $\sum_{j=k}^{m-1}\binom jk = \binom m{k+1}$، وهي
متطابقة عصا الهوكي.

\textbf{23.} تعطي جداول الفروق (أو النشر المباشر)
\[
X^2 = B_1 + 2 B_2, \qquad X^3 = B_1 + 6 B_2 + 6 B_3
\]
(وللتحقق: $B_1 + 2B_2 = X + X(X-1) = X^2$؛ وعند $X = 1, 2, 3$ يعطي
الثاني $1, 8, 27$). ويعطي السؤال 22 عندئذ
\[
\sum_{j=0}^{m-1} j^2 = B_2(m) + 2B_3(m)
= \binom m2 + 2\binom m3 = \frac{m(m-1)(2m-1)}{6},
\]
\[
\sum_{j=0}^{m-1} j^3 = B_2(m) + 6B_3(m) + 6B_4(m)
= \binom m2 + 6\binom m3 + 6\binom m4 .
\]
وبنشر المقدار الأخير: $\binom m2 + 6\binom m3 + 6\binom
m4 = \frac{m(m-1)}{2}\bigl[1 + 2(m-2) +
\frac{(m-2)(m-3)}{2}\bigr] = \frac{m^2(m-1)^2}{4} = \binom m2^2$.
وباستبدال $m + 1$ بالمقدار $m$: $1^3 + \dots + m^3 =
\bigl(\frac{m(m+1)}2\bigr)^2 = (1 + \dots + m)^2$، وهي متطابقة
نيقوماخوس.

\textbf{24.} العقد $0, 1, 2$، \omterm{def:b1:poly:def}{وكثير الحدود} $X^2$. \omterm{def:b1:vspaces:free}{الأساس} الوحيد الحدّ
$(1, X, X^2)$: \omterm{prop:b1:vspaces:coordinates}{الإحداثيات} $(0, 0, 1)$. \omterm{def:b1:vspaces:free}{وأساس} لاغرانج:
\omterm{prop:b1:vspaces:coordinates}{الإحداثيات} هي القيم $(0, 1, 4)$ (السؤال 4). \omterm{def:b1:vspaces:free}{وأساس} نيوتن
$(1, X, X(X-1))$: \omterm{pb:b1:vspaces:1}{الفروق المقسومة} $f[0] = 0$ و $f[0,1] = 1$ و
$f[0,1,2] = \frac{3 - 1}{2} = 1$ (السؤال 9)، ومنه \omterm{prop:b1:vspaces:coordinates}{الإحداثيات}
$(0, 1, 1)$ --- وفعلًا $X + X(X-1) = X^2$. أي ثلاثة أسس وثلاثة
متجهات \omterm{prop:b1:vspaces:coordinates}{إحداثيات} \omterm{def:b1:poly:def}{وكثير حدود} واحد.

\textbf{25.} (أ) السؤال 4 تلقائي لأن $(L_i)$
\emph{\omterm{def:b1:vspaces:free}{أساس}}: فوجود المقايسة ووحدانيتها هما
بالضبط وجود \omterm{prop:b1:vspaces:coordinates}{الإحداثيات} ووحدانيتها. (ب) \omterm{def:b1:vspaces:free}{وأساس} نيوتن
سلّم، ومنه تُحسب \omterm{prop:b1:vspaces:coordinates}{الإحداثيات} بقسمات متتالية ---
إذ تضيف كل عقدة جديدة حدًّا واحدًا دون إزعاج
السابقة --- بينما \omterm{prop:b1:vspaces:coordinates}{إحداثيات} $P$ في \omterm{def:b1:vspaces:free}{أساس} لاغرانج هي
القيم $P(x_i)$، متوفرة دون أيّ حساب أصلًا. (ج) والأساسان
كلاهما حرّ بالمحكين نفسيهما في
\cref{prop:b1:vspaces:freecriteria}: الدرجات المتمايزة من أجل نيوتن،
والتقويم عند العقد من أجل لاغرانج. (د) وتقول مبرهنة بوليا
إن «$P(\Z) \subseteq \Z$»، وهي خاصية للقيم، تكافئ
صحيحية \omterm{prop:b1:vspaces:coordinates}{الإحداثيات} في \omterm{def:b1:vspaces:free}{الأساس} $(B_k)$ --- فحساب
\omterm{def:b1:poly:def}{كثير الحدود} لا يصير مرئيًا إلا في \omterm{def:b1:vspaces:free}{الأساس}
المكيَّف مع السؤال.
\end{solution}
